% GENERATED FILE. Edit canonical lessons, then run npm run generate:books.
% canonical-source-hash: fnv1a64:28c2233ed3142aaa
% canonical-lessons: RU-C168-svadba, RU-C168-zhenikh, RU-C168-nevesta, RU-C168-rodstvennik, RU-C168-pensioner

\chapter{Wedding, Bridegroom, Bride, Relative, Pensioner}
\label{ch:ru-wedding-bridegroom-bride-relative-pensioner}
\hlchaptermodality{168}

\begin{chapteropening}
\textbf{By the end of this chapter:} \emph{I can say a wedding, a bridegroom, a bride, a relative and a pensioner.}

Complete the last lesson of chapter 168: I can say a wedding, a bridegroom, a bride, a relative and a pensioner.
\end{chapteropening}

\section[\texorpdfstring{\ru{свадьба} (svádba)}{свадьба (svádba)}]{\ru{свадьба} (svádba) — a wedding}

\label{lesson:RU-C168-svadba}



\begin{quote}
Before the new one: say the Russian for petrol, then the Russian for a petrol station.
\end{quote}

\subsection*{You'll want to know: \ru{свадьба}}
\textbf{\ru{свадьба}} (\emph{svádba}) — \textquotedblleft{}a wedding\textquotedblright{}.

\begin{cousinweb}
\textbf{\ru{свадьба}} — Old Russian \emph{\ru{сватьба}}; it comes from \emph{\ru{сват}}, \textquotedblleft{}a matchmaker, an in-law\textquotedblright{} (Vasmer), and first meant \textquotedblleft{}match-making\textquotedblright{}. Vasmer connects \emph{\ru{сват}} with \emph{\ru{свой}}, \textquotedblleft{}one's own\textquotedblright{}.
\end{cousinweb}

\begin{grammarlens}[title={What you've built}]
One more word for this chapter.
\end{grammarlens}

\subsection*{Your turn}
\begin{itemize}
\raggedright
  \item \emph{Say it:} \emph{svádba}
  \item \emph{Say it:} \emph{svádba}, once more
  \item \emph{Say it:} say \emph{svádba}
  \item \emph{Recall:} say the Russian for petrol, then the Russian for a petrol station, then say \emph{svádba} again
\end{itemize}

\begin{tcolorbox}[breakable,colback=teal!4,colframe=teal!35!black,title={Before you move on}]
What does \ru{свадьба} mean? (\textquotedblleft{}A wedding\textquotedblright{}.) Say it once more.
\end{tcolorbox}
\section[\texorpdfstring{\ru{жених} (zheníkh)}{жених (zheníkh)}]{\ru{жених} (zheníkh) — a bridegroom, a fiancé}

\label{lesson:RU-C168-zhenikh}



\begin{quote}
Before the new one: say the Russian for a petrol station, then the Russian for a wedding.
\end{quote}

\subsection*{You'll want to know: \ru{жених}}
\textbf{\ru{жених}} (\emph{zheníkh}) — \textquotedblleft{}a bridegroom, a fiancé\textquotedblright{}.

\begin{cousinweb}
\textbf{\ru{жених}} is built on \emph{\ru{женить}}, \textquotedblleft{}to marry (a man) off\textquotedblright{}, from \emph{\ru{жена}}, \textquotedblleft{}wife\textquotedblright{} (Vasmer).
\end{cousinweb}

\begin{grammarlens}[title={What you've built}]
One more word for this chapter.
\end{grammarlens}

\subsection*{Your turn}
\begin{itemize}
\raggedright
  \item \emph{Say it:} \emph{zheníkh}
  \item \emph{Say it:} \emph{zheníkh}, once more
  \item \emph{Say it:} say \emph{zheníkh}
  \item \emph{Recall:} say the Russian for a petrol station, then the Russian for a wedding, then say \emph{zheníkh} again
\end{itemize}

\begin{tcolorbox}[breakable,colback=teal!4,colframe=teal!35!black,title={Before you move on}]
What does \ru{жених} mean? (\textquotedblleft{}A bridegroom, a fiancé\textquotedblright{}.) Say it once more.
\end{tcolorbox}
\section[\texorpdfstring{\ru{невеста} (nevésta)}{невеста (nevésta)}]{\ru{невеста} (nevésta) — a bride, a fiancée}

\label{lesson:RU-C168-nevesta}



\begin{quote}
Before the new one: say the Russian for a wedding, then the Russian for a bridegroom.
\end{quote}

\subsection*{You'll want to know: \ru{невеста}}
\textbf{\ru{невеста}} (\emph{nevésta}) — \textquotedblleft{}a bride, a fiancée\textquotedblright{}.

\begin{cousinweb}
\textbf{\ru{невеста}} — Vasmer still finds the old explanation best: \textquotedblleft{}the unknown one\textquotedblright{}, \emph{\ru{не}}, \textquotedblleft{}not\textquotedblright{}, with the root of \emph{\ru{ведать}}, \textquotedblleft{}to know\textquotedblright{}, for the new bride is a stranger in her husband's house. Trubachev's note cites other views.
\end{cousinweb}

\begin{grammarlens}[title={What you've built}]
One more word for this chapter.
\end{grammarlens}

\subsection*{Your turn}
\begin{itemize}
\raggedright
  \item \emph{Say it:} \emph{nevésta}
  \item \emph{Say it:} \emph{nevésta}, once more
  \item \emph{Say it:} say \emph{nevésta}
  \item \emph{Recall:} say the Russian for a wedding, then the Russian for a bridegroom, then say \emph{nevésta} again
\end{itemize}

\begin{tcolorbox}[breakable,colback=teal!4,colframe=teal!35!black,title={Before you move on}]
What does \ru{невеста} mean? (\textquotedblleft{}A bride, a fiancée\textquotedblright{}.) Say it once more.
\end{tcolorbox}
\section[\texorpdfstring{\ru{родственник} (ródstvennik)}{родственник (ródstvennik)}]{\ru{родственник} (ródstvennik) — a relative}

\label{lesson:RU-C168-rodstvennik}



\begin{quote}
Before the new one: say the Russian for a bridegroom, then the Russian for a bride.
\end{quote}

\subsection*{You'll want to know: \ru{родственник}}
\textbf{\ru{родственник}} (\emph{ródstvennik}) — \textquotedblleft{}a relative\textquotedblright{}.

\begin{cousinweb}
\textbf{\ru{родственник}} is built on \emph{\ru{родство}}, \textquotedblleft{}kinship\textquotedblright{}, from \emph{\ru{род}}, \textquotedblleft{}family, kin\textquotedblright{}, with \emph{-\ru{ник}}. Vasmer relates \emph{\ru{род}} to Latvian \emph{rads}, \textquotedblleft{}a relative\textquotedblright{}, and Sanskrit \emph{várdhati}, \textquotedblleft{}grows\textquotedblright{}.
\end{cousinweb}

\begin{grammarlens}[title={What you've built}]
One more word for this chapter.
\end{grammarlens}

\subsection*{Your turn}
\begin{itemize}
\raggedright
  \item \emph{Say it:} \emph{ródstvennik}
  \item \emph{Say it:} \emph{ródstvennik}, once more
  \item \emph{Say it:} say \emph{ródstvennik}
  \item \emph{Recall:} say the Russian for a bridegroom, then the Russian for a bride, then say \emph{ródstvennik} again
\end{itemize}

\begin{tcolorbox}[breakable,colback=teal!4,colframe=teal!35!black,title={Before you move on}]
What does \ru{родственник} mean? (\textquotedblleft{}A relative\textquotedblright{}.) Say it once more.
\end{tcolorbox}
\section[\texorpdfstring{\ru{пенсионер} (pensionér)}{пенсионер (pensionér)}]{\ru{пенсионер} (pensionér) — a pensioner}

\label{lesson:RU-C168-pensioner}



\begin{quote}
Before the new one: say the Russian for a bride, then the Russian for a relative.
\end{quote}

\subsection*{You'll want to know: \ru{пенсионер}}
\textbf{\ru{пенсионер}} (\emph{pensionér}) — \textquotedblleft{}a pensioner\textquotedblright{}.

\begin{cousinweb}
\textbf{\ru{пенсионер}} comes from French \emph{pensionnaire}, built on \emph{pension}, from Latin \emph{pēnsiō}, \textquotedblleft{}a payment\textquotedblright{}, from \emph{pendere}, \textquotedblleft{}to weigh, to pay\textquotedblright{}.
\end{cousinweb}

\begin{grammarlens}[title={What you've built}]
One more word for this chapter.
\end{grammarlens}

\subsection*{Your turn}
\begin{itemize}
\raggedright
  \item \emph{Say it:} \emph{pensionér}
  \item \emph{Say it:} \emph{pensionér}, once more
  \item \emph{Say it:} say \emph{pensionér}
  \item \emph{Recall:} say the Russian for a bride, then the Russian for a relative, then say \emph{pensionér} again
\end{itemize}

\begin{tcolorbox}[breakable,colback=teal!4,colframe=teal!35!black,title={Before you move on}]
What does \ru{пенсионер} mean? (\textquotedblleft{}A pensioner\textquotedblright{}.) Say it once more.
\end{tcolorbox}
